\bookchapter{Глава 58. Справочник и ресурсы}

\chapterintro{Эта глава — справочник, которым можно пользоваться и отдельно. Когда не можете вспомнить, как что-то пишется, или нужно быстро принять решение, — просто заглядывайте сюда. Советуем обращаться с ним как с «карманными карточками»: шпаргалки по часто используемому коду, частые вопросы, учебные ресурсы — ни одна страница не пропадёт.}

\booksection{Шпаргалка по часто используемому коду}

\begin{center}
\begin{tabularx}{\linewidth}{|X|X|}
\hline
\textbf{Что нужно сделать} & \textbf{Ключевой код} \tabularnewline
\hline
Запомнить изменяемое состояние & var x by remember \{ mutableStateOf(начальное Значение) \} \tabularnewline
Не терять состояние при повороте экрана & rememberSaveable \{ mutableStateOf(начальное Значение) \} \tabularnewline
Вертикальная/горизонтальная компоновка & \texttt{Column \{ \}}/\texttt{Row \{ \}} \tabularnewline
Список с ленивой загрузкой & LazyColumn \{ items(список) \{ ... \} \} \tabularnewline
Переход между экранами & navController.navigate("маршрут") \tabularnewline
Долгая работа в корутине & \texttt{viewModelScope.\allowbreak launch \{ .\allowbreak .\allowbreak .\allowbreak  \}} \tabularnewline
Наблюдать за потоком данных & \texttt{val list by flow.\allowbreak collectAsState()} \tabularnewline
Хранить небольшие настройки & context.dataStore.edit \{ it[KEY] = значение \} \tabularnewline
Хранить структурированные данные & Тройка Room: @Entity / @Dao / @Database \tabularnewline
Получать данные из сети & Интерфейс Retrofit + suspend-функции \tabularnewline
Отправлять уведомления & NotificationChannel + NotificationCompat.Builder \tabularnewline
Писать в лог & \texttt{Log.\allowbreak d("TAG", "message")} \tabularnewline
\hline
\end{tabularx}
\end{center}

\begin{center}\small\color{gray}Шпаргалка по разработке под Android (версия августа 2026 года)\end{center}

\booksection{Частые вопросы (FAQ)}

В 1: Я совсем не умею программировать. Сколько времени займёт освоение разработки под Android?

Идите по пути этой книги: недели 1–2 — установка окружения и основы Kotlin; недели 3–4 — изучение Compose и создание интерфейса; недели 5–6 — завершение проекта «учёт расходов»; недели 7–8 — завершение проекта «погода». Два месяца — от нуля до умения самостоятельно делать небольшое приложение. При условии, что вы готовы уделять 1–2 часа в день и писать код руками, а не только читать.

В 2: Нужно ли ещё учить Java?

Нет. Новые проекты, новые учебники и новые вакансии почти все ориентированы на Kotlin. Много Java вы встретите только при «поддержке проекта десятилетней давности». Если после Kotlin работа потребует Java, подтянуть его получится быстро — в основе обоих языков одна экосистема виртуальной машины.

В 3: Эмулятор или реальное устройство?

Для отладки в процессе разработки удобен эмулятор (снимки состояния, много моделей устройств), но финальная приёмка — обязательно на реальном устройстве: ощущения от касаний, производительность, настоящая сеть, датчики — эмулятор этого не заменит. Если есть возможность, купить реальное устройство среднего класса под тесты — самая выгодная инвестиция.

В 4: Легко ли найти работу в разработке под Android?

Ситуация 2026 года: на начальные позиции уровня «просто умею делать интерфейсы» конкуренция острая, но разработчики, которые «знают Kotlin + понимают архитектуру + имеют опубликованное приложение», по-прежнему нарасхват. Лучшее резюме для вас, пока вы учитесь, — это сделанное вами приложение в магазине.

В 5: Что делать, если застряли на ошибке?

Золотые три шага: (1) скопируйте первую строку ошибки как есть и поищите в StackOverflow/поисковике; (2) посмотрите соответствующую страницу официальной документации; (3) если не помогло, соберите «ожидаемое поведение + фактический результат + полный текст ошибки + относящийся код» и отправьте вопрос ИИ или сообществу. Подавляющее большинство проблем решается уже на первом шаге.

\booksection{Рекомендуемые учебные ресурсы}

\begin{center}
\begin{tabularx}{\linewidth}{|X|X|X|}
\hline
\textbf{Тип} & \textbf{Ресурс} & \textbf{Почему стоит попробовать} \tabularnewline
\hline
Официальный & developer.android.com & Официальная документация Android: внутри и руководства по Compose/Kotlin — источник номер один \tabularnewline
Официальный & Официальный сайт Kotlin (kotlinlang.org) & Официальная документация языка + онлайн-песочница \tabularnewline
Официальный & Google Codelabs & Пошаговые практические уроки «рука об руку»; пятёрка по дружелюбности к новичкам \tabularnewline
Сообщество & Stack Overflow (метка android) & Первая остановка за помощью при ошибках \tabularnewline
Сообщество & Juejin / раздел Android на CSDN & Много практических статей на китайском; проверяйте свежесть (сначала смотрите дату) \tabularnewline
Видео & Официальный канал Google Developers & Разборы новых технологий с ежегодных конференций разработчиков \tabularnewline
Из этой книги & Часть 7 «Практика программирования с ИИ» (главы 37–47) & Полная методика разработки под Android с помощью ИИ \tabularnewline
\hline
\end{tabularx}
\end{center}

\begin{center}\small\color{gray}Список учебных ресурсов по разработке под Android (август 2026 года)\end{center}

\booksection{Шпаргалка по командам Gradle}

В обычной жизни достаточно нажимать кнопки, но командная строка иногда эффективнее (особенно в связке с ИИ-инструментами или CI-конвейером):

\begin{center}
\begin{tabularx}{\linewidth}{|X|X|}
\hline
\textbf{Команда} & \textbf{Назначение} \tabularnewline
\hline
\texttt{.\allowbreak /\allowbreak gradlew assembleDebug} & Собрать отладочный APK (в Windows уберите ./) \tabularnewline
\texttt{.\allowbreak /\allowbreak gradlew assembleRelease} & Собрать релизную версию (нужна настроенная подпись) \tabularnewline
\texttt{.\allowbreak /\allowbreak gradlew test} & Запустить все модульные тесты \tabularnewline
\texttt{.\allowbreak /\allowbreak gradlew lint} & Статический анализ: находит потенциальные баги и неаккуратный код \tabularnewline
\texttt{.\allowbreak /\allowbreak gradlew clean} & Очистить результаты сборки («перезагрузка» в мире компиляции) \tabularnewline
\texttt{.\allowbreak /\allowbreak gradlew :app: dependencies} & Показать дерево зависимостей (волшебное средство против конфликтов версий) \tabularnewline
\hline
\end{tabularx}
\end{center}

\begin{center}\small\color{gray}Шпаргалка по командам Gradle\end{center}

\begin{tcbinfo}{Совет: если сборка барахлит, сначала clean}
Если «вчера всё работало, а сегодня нет», то в шести-семи случаях из десяти виноват перепутанный кэш сборки. Выполните ./gradlew clean и соберите заново — это лечит большую часть «мистических» проблем. Не помогло — перезапустите Android Studio; этот приём лучше всего работает в паре с File Invalidate Caches.

\end{tcbinfo}

\booksection{Заключение: от первого приложения до N-го}

\booksection{Заключение: от первого приложения до N-го}

Оглядываясь назад: вы прошли путь от «не знаю, что такое Activity» до установленного окружения, освоенного Kotlin, интерфейсов, нарисованных в Compose, подключённых к приложению базы данных и сети, а дальше — тестирования, сборки пакета и публикации. Самый трудный отрезок пути «от нуля к единице» вы уже прошли.

Мир технологий меняется быстро: фреймворки обновляются, инструменты сменяют друг друга, а ИИ берёт на себя всё больше «клавиатурной» работы. Но одна вещь не устареет никогда: ваша способность превратить расплывчатую идею в чёткие шаги и шаг за шагом её воплотить. Инструменты — усилитель, образ мышления — фундамент. Усилитель можно заменить, фундамент шататься не должен.

Вы больше не просто «умеющий пользоваться телефоном». Вы можете превратить идеи из головы в настоящие приложения, работающие на трёхстах миллионах Android-устройств. С сегодняшнего дня творите с клавиатурой — в этом романтика Android-разработчика.

\booksection{Список литературы}

\begin{enumerate}
  \item Android Developers. (2026). Android Developer Guides and Jetpack Compose Documentation. developer.android.com.
  \item Kotlin Foundation. (2026). Kotlin Language Documentation. kotlinlang.org.
  \item Android Open Source Project. (2026). Android Platform Version History and API Levels.
  \item Google Play. (2026). Play Console Help: Publishing and App Signing.
  \item AndroidX Release Notes. (2026). Room / DataStore / Navigation / WorkManager releases.
  \item Square. (2026). Retrofit Documentation. github.com/square/retrofit.
\end{enumerate}
